\section{Conclusion}
\label{sec:conclusion}

Suitability scores conflate two things a decision-maker must separate: how hard a
condition is, and how skilled the person facing it is. We showed that a
continuous-item Item Response Theory model---$P(y{=}1) = \sigmoid(a(\skill_\rider -
\diff(\metric,\site)))$ with a smooth, physics-anchored difficulty function---
identifies both from binary behavioural outcomes alone. The model is strictly more
general than a single suitability curve, which it recovers exactly on marginalising
skill; it is identified when the rider${\times}$condition incidence graph is
connected, and it fails safely to the single curve when it is not. A synthetic
recovery study confirms the estimator recovers latent skill at correlation $0.96$,
locates the difficulty minimum within $\pm 3$ units of ground truth, and improves
held-out Brier Skill Score by $+0.33$ over the expert-curve baseline---all
reproducible from one command on synthetic data, with no proprietary observations.

Beyond the estimator, the recovered difficulty function defines a measurable,
anonymous, site-level construct---an intrinsic difficulty atlas---that existing
meteorological networks do not capture and that is joinable to them. Establishing
this construct, and the conditions under which it can be identified and governed,
is the paper's lasting contribution. What remains is empirical: to fit the model
to real cohorts as they accumulate and test whether human behaviour in the field
follows the structure that, in simulation, it recovers so cleanly.
